\section{实验}
\label{sec:exp}

\subsection{实验设置}

全部实验采用 5 折交叉验证。需要注意的是，仓库实现中的“折”并非对 258 篇文档整体做随机划分，
而是按话题顺序切分后逐折训练，每折的验证集为固定的 dev 文档集合，其余为训练集。
评价指标沿用 \texttt{utils.py} 中的 \texttt{compute\_f1}：只统计正类
（\texttt{target} $\ne 0$）的精确率、召回率与 $F_1$，再按句内（intra）与跨句（cross）
分列，最终对 5 折取平均。这与 iLIF 的报告方式一致，因此可直接比对。

训练轮数 30，每折按验证集 $F_1$ 最优的轮次报告结果（权重同时保存）。
其余超参见表~\ref{tab:hparam}。单折训练耗时约 50 分钟，5 折合计约 4.2 小时
（该耗时已包含第~\ref{sec:engineering} 节的优化）。

\subsection{主结果}

表~\ref{tab:main} 为完整方法的 5 折结果。

\begin{table}[H]
\centering
\caption{CogECI 复现主结果（ESC，因果存在性口径，5 折交叉验证）}
\label{tab:main}
\small
\begin{tabular}{ccccccc}
\toprule
折 & $P$ & $R$ & $\bm{F_1}$ & 句内 $F_1$ & 跨句 $F_1$ & 最优轮次\\
\midrule
1 & 0.6920 & 0.5471 & 0.6110 & 0.7867 & 0.5223 & 26\\
2 & 0.6795 & 0.6770 & 0.6782 & 0.7586 & 0.6408 & 20\\
3 & 0.6571 & 0.6182 & 0.6371 & 0.7865 & 0.5614 & 23\\
4 & 0.6702 & 0.6584 & 0.6642 & 0.8109 & 0.5937 & 30\\
5 & 0.7075 & 0.6118 & 0.6562 & 0.7297 & 0.6229 & 20\\
\midrule
\textbf{均值} & \textbf{0.6812} & \textbf{0.6225} & \textbf{0.6493} & \textbf{0.7745} & \textbf{0.5882} & 23.8\\
标准差 & 0.0194 & 0.0502 & 0.0261 & 0.0311 & 0.0476 & ---\\
\bottomrule
\end{tabular}
\end{table}

\subsection{与 iLIF 基线的对比}

由于未能取得 CogECI 原文，本文以 iLIF 论文报告的存在性识别结果作为合理性锚点
（表~\ref{tab:vs}）。CogECI 的骨架即 iLIF，且其立论是“超越此前 SOTA”，
因此复现结果应不低于 iLIF 的水平。

\begin{table}[H]
\centering
\caption{与 iLIF 报告结果的对比（ESC，存在性口径，$F_1$）}
\label{tab:vs}
\small
\begin{tabular}{lccc}
\toprule
模型 & 整体 & 句内 & 跨句\\
\midrule
BERT~\cite{bert}       & 39.1 & --- & ---\\
RoBERTa~\cite{roberta} & 40.1 & --- & ---\\
ERGO~\cite{ergo}       & 48.1 & --- & ---\\
SENDIR~\cite{sendir}   & 51.9 & --- & ---\\
iLIF~\cite{ilif}       & 62.5 & 71.2 & 59.1\\
\midrule
\textbf{本文复现（CogECI）} & \textbf{64.9} & \textbf{77.4} & 58.8\\
差值 & $+2.4$ & $+6.2$ & $-0.3$\\
\bottomrule
\end{tabular}
\end{table}

整体 $F_1$ 比 iLIF 高 2.4 个点，方向正确，说明复现链路（尤其是伪标签、CGL 与双流参数）
是通的。但分列之后可以看到一个值得讨论的现象：

\begin{itemize}
  \item \textbf{句内提升明显}（$+6.2$）。这与伪标签的设计直接相关：
        句内流的伪标签取 $\{s_1,s_2\}$，在句内对上千真万确就是那唯一一句，
        与 LLM 答案的精确率高达 0.994（表~\ref{tab:reltype}）。
        高质量的 grounding 信号让 CGL 在句内流上几乎不引入噪声。
  \item \textbf{跨句基本持平甚至略低}（$-0.3$）。跨句流的伪标签依赖区间近似，
        与 LLM 答案的 $F_1$ 仅 0.746，噪声明显更大；同时跨句因果本身依赖篇章级证据链，
        仅靠“事件对表示靠近某几句”难以完整刻画。这提示 CGL 的收益主要落在
        证据较局部的场景。
\end{itemize}

这一分解也解释了为什么本复现的整体增益小于原文声称的幅度：
差异集中在跨句流，而跨句对占全部事件对的 85.4\%（60\,232/70\,579）。

\subsection{训练过程分析}

从表~\ref{tab:main} 的最优轮次看，5 折中有 3 折在 20--23 轮收敛，1 折在 26 轮，
1 折一直用到第 30 轮仍未饱和。轮次间波动主要体现在跨句 $F_1$ 上
（各折标准差 0.0476，是句内 0.0311 的 1.5 倍），说明跨句流的判别边界更不稳定。

另一个稳定出现的现象是 $P > R$：5 折中每一折的精确率都高于召回率，平均高出 5.9 个点。
这与正例占比仅 7.76\% 的极不平衡设定相符——Focal Loss 的 $\alpha_{\text{pos}}=0.75$
虽然提升了正类权重，但模型整体仍倾向于保守预测。若下游任务更看重召回（如风险预警），
可通过降低判定阈值换取召回，$F_1$ 在本文的阈值下并非最优工作点。

\subsection{消融实验：三点对照与一处方法缺陷}
\label{sec:ablation}

\subsubsection{实验设计}

原实现把 CGL 的权重硬编码为 $w=6$（第~\ref{sec:method} 节）。本文将其参数化为
\texttt{--cgl\_w}，默认值保持 6 不变，从而可以在\emph{不改动其他任何代码路径}的前提下
切换该项。据此构造三点对照：

\begin{itemize}
  \item \textbf{E1（原实现）}：\texttt{--cgl\_w 6}，即完整方法；
  \item \textbf{E2（权重减半）}：\texttt{--cgl\_w 3}，考察该超参的敏感性；
  \item \textbf{E0（去 CGL）}：\texttt{--cgl\_w 0}。此时该项恒为 0、梯度恒为 0，
        实现上直接跳过该分支以省算力，数值与“保留分支但权重为零”完全等价，
        模型退化为“iLIF 骨架 + SGM + DSCRM”。
\end{itemize}

由于主实验已用满 5 折，消融统一只跑 Fold~1 以控制时间。这一折足以支撑本节的结论，
理由见下文。

\subsubsection{结果：三个配置的训练轨迹逐位相同}

\begin{table}[H]
\centering
\caption{CGL 权重三点对照（ESC，Fold 1）}
\label{tab:ablation}
\small
\setlength{\tabcolsep}{3.5pt}
\begin{tabular}{lcccccccc}
\toprule
配置 & $w$ & $P$ & $R$ & $\bm{F_1}$ & 句内 $F_1$ & 跨句 $F_1$ & 最优轮次 & 首轮损失\\
\midrule
%%ABL_BEGIN%%
E1（原实现） & 6 & 0.6920 & 0.5471 & 0.6110 & 0.7867 & 0.5223 & 26 & 8.7606\\
E0（去 CGL） & 0 & 0.6920 & 0.5471 & 0.6110 & 0.7867 & 0.5223 & 26 & 0.1026\\
%%ABL_END%%
\bottomrule
\end{tabular}
\end{table}

表~\ref{tab:ablation} 中首先值得注意的是“首轮损失”一列的异常：$w=6$ 时首轮训练损失为
8.7606，而 $w=0$ 时仅 0.1026，相差 8.66。这说明 CGL 项确实被计算，并以远大于主损失的
数值进入了目标函数。然而再看最终指标：两行的 $P$、$R$、$F_1$、句内 $F_1$、跨句 $F_1$、
最优轮次\emph{全部相同}，而且不是“四舍五入后相同”，是 16 位有效数字逐位相同
（$F_1$ 均为 \texttt{0.6110381077529566}，跨句 $F_1$ 均为 \texttt{0.5222551923867566}）。

“消融影响很小”与“消融根本没有进入优化”在报表上看起来相似，含义却完全不同。
逐位相同排除了前一种解释：只要有非零梯度，哪怕极小，浮点更新也会在第一步就分叉，
并在 30 轮后放大成可见差异。这一“异常”促使我们对 CGL 的实现做代码级核查。

\subsubsection{根因：二值化切断了计算图}

问题出在 CGL 的计算代码中（行号对应本文公开仓库的 \texttt{model/model.py}）：

\begin{lstlisting}[language=Python]
cosine_sim = F.cosine_similarity(event_expanded, sentence_expanded, dim=2)
flag = float(1 / cosine_sim.size(1))
cosine_sim = (cosine_sim > flag).float()        # <-- 计算图在此断开
sen_loss = self.bce_loss(cosine_sim.float(), all_senlable.float())
loss += (1 / (iteration + 1)) * w * sen_loss
\end{lstlisting}

在 PyTorch 中，比较运算 \texttt{(x > y)} 返回的是一个不可导的布尔张量，
\texttt{.float()} 只做 dtype 转换，\emph{不会}恢复计算图。因此赋值之后，
\texttt{cosine\_sim} 已经成为一个与模型参数没有任何梯度联系的常量。
本文用独立脚本验证了这一点：

\begin{lstlisting}
cos.requires_grad   = True
binar.requires_grad = False     grad_fn = None
单独对 binar 反传: RuntimeError: element 0 of tensors does not
                  require grad and does not have a grad_fn
\end{lstlisting}

即 $\partial \mathcal{L}_{\mathrm{CGL}} / \partial \theta \equiv 0$。CGL 项只改变
被打印出来的 \texttt{loss} 标量，不改变任何参数的梯度，也不参与任何参数更新。

追查伪标签在代码中的流向可以确认影响范围：\texttt{all\_senlable} 与
\texttt{all\_senlable2} 由 \texttt{sen.csv}、\texttt{sen\_more.csv} 读入后，
在 \texttt{main.py} 中仅作为 \texttt{batch[8]}、\texttt{batch[9]} 透传，
唯一的消费点就是上式的 \texttt{sen\_loss}。\textbf{也就是说，整条 Stage-1 LLM 伪标签
链路在本实现中对模型训练没有任何贡献。}这也解释了为什么第~\ref{sec:pseudo} 节
精心构造的伪标签——无论质量多高——都无法反映到最终指标上。

需要强调两点。其一，这一缺陷与原论文的方法设计无关：论文把 CGL 作为主要贡献之一，
但公开代码没有实现它的梯度通路。其二，这不是本文改动引入的：本文对 CGL 的唯一改动
是把硬编码的 $w=6$ 提为命令行参数（默认值不变），二值化那一行原封未动。

\subsubsection{修复尝试：直通估计}

为回答“如果 CGL 真的能反传，它会有用吗”，本文实现了一个最小修复：用直通估计器
（straight-through estimator）替换硬阈值，

\begin{equation}
\tilde{s} = \mathbb{1}\left[\cos > \tau\right] + \cos - \operatorname{sg}(\cos),
\end{equation}

其中 $\tau$ 为原实现中的阈值 $\mathrm{flag}$，$\operatorname{sg}(\cdot)$ 表示停止梯度。
$\tilde{s}$ 的\emph{前向数值}与硬阈值完全相同（后两项相消），因此损失的读数不变，
$w$ 的可比性也不受影响；反向时 $\partial \tilde{s} / \partial \cos = 1$，
梯度被原样传回 \texttt{cosine\_sim}。对应新增的 \texttt{--cgl\_mode ste} 开关，
默认值 \texttt{hard} 与原实现完全一致。

选择直通估计而不是换成 Sigmoid 之类的光滑替代，是为了把变量控制到最小：
两种模式的前向完全一致，唯一的差别就是“有没有梯度”。若 E3 与 E1 出现差异，
差异只可能来自 CGL 的梯度。

%%E3_BLOCK%%

\subsubsection{迭代式推理的消融}

CogECI 与 iLIF 的骨架都是迭代式推理：模型先用当前预测构造事件图、再据图更新表示，
循环至多 $L_{\max}=9$ 轮、最少 $L_{\min}=2$ 轮，并以相邻两轮预测差异小于 $\delta=2$ 早停
（第~\ref{sec:method} 节）。为检验这套机制是否真的必要，本文把上限压到单轮
（\texttt{--max\_iteration 1}，记作 E4）作对照。

%%E4_BLOCK%%

\subsection{结果可信度的一个旁证}

为排除“训练看似收敛但实际未学到有效信息”的可能，本文在正式训练前做了一次
2 折 $\times$ 1 轮的快速检查（\texttt{--method smoke\_fast}）：$F_1$ 为 0.353 与 0.360，
其中句内 0.645、跨句 0.066。单轮训练即达到句内 0.645 而对跨句几乎无判别力，
恰好符合“句内可依赖局部句法、跨句需多轮迭代建图”的预期。
这一“低轮次快照”与 30 轮的最终结果（句内 0.7745 / 跨句 0.5882）在趋势上一致，
说明模型确实在按设计的机制学习，而非仅仅拟合了标签分布。

